% Opis techniczny projektu Email Triage Agent
% Kompilacja: Overleaf, kompilator pdfLaTeX (domyślny)

\documentclass[11pt,a4paper]{article}

\usepackage[utf8]{inputenc}
\usepackage[T1]{fontenc}
\usepackage[polish]{babel}
\usepackage{lmodern}
\usepackage[a4paper,margin=2.5cm]{geometry}
\usepackage{microtype}
\usepackage{xcolor}
\usepackage{booktabs}
\usepackage{tabularx}
\usepackage{array}
\usepackage{enumitem}
\usepackage{listings}
\usepackage{tikz}
\usetikzlibrary{positioning,arrows.meta,shapes.geometric,fit,babel}
\usepackage{hyperref}

\hypersetup{
  colorlinks=true,
  linkcolor=blue!50!black,
  urlcolor=blue!50!black,
  pdftitle={Email Triage Agent - opis techniczny},
}

\setlength{\parindent}{0pt}
\setlength{\parskip}{0.6em}
\setlist{itemsep=0.2em}

% --- kolory i styl listingów ---
\definecolor{codebg}{RGB}{246,246,246}
\definecolor{codekw}{RGB}{0,70,140}
\definecolor{codestr}{RGB}{150,40,40}
\definecolor{codecom}{RGB}{90,110,90}

\lstset{
  language=Python,
  basicstyle=\ttfamily\small,
  keywordstyle=\color{codekw}\bfseries,
  stringstyle=\color{codestr},
  commentstyle=\color{codecom}\itshape,
  backgroundcolor=\color{codebg},
  frame=single,
  rulecolor=\color{black!20},
  breaklines=true,
  columns=fullflexible,
  keepspaces=true,
  showstringspaces=false,
  xleftmargin=0.5em,
  xrightmargin=0.5em,
  aboveskip=0.8em,
  belowskip=0.8em,
  extendedchars=true,
  literate=
    {ą}{{\k{a}}}1 {ć}{{\'c}}1 {ę}{{\k{e}}}1 {ł}{{\l{}}}1 {ń}{{\'n}}1
    {ó}{{\'o}}1 {ś}{{\'s}}1 {ź}{{\'z}}1 {ż}{{\.z}}1
    {Ą}{{\k{A}}}1 {Ć}{{\'C}}1 {Ę}{{\k{E}}}1 {Ł}{{\L{}}}1 {Ń}{{\'N}}1
    {Ó}{{\'O}}1 {Ś}{{\'S}}1 {Ź}{{\'Z}}1 {Ż}{{\.Z}}1,
}

\newcommand{\plik}[1]{\texttt{#1}}

% --- style diagramów ---
\tikzset{
  wezel/.style={rectangle, rounded corners=3pt, draw=black!70, fill=blue!8,
    minimum width=2.8cm, minimum height=0.9cm, align=center, font=\small},
  llm/.style={wezel, fill=orange!15},
  koniec/.style={wezel, fill=black!8},
  start/.style={ellipse, draw=black!70, fill=green!10, minimum height=0.8cm,
    align=center, font=\small},
  zew/.style={rectangle, draw=black!60, dashed, fill=yellow!10, align=center,
    font=\small, minimum height=0.9cm},
  strz/.style={-{Stealth[length=2.2mm]}, thick, black!70},
  etykieta/.style={font=\footnotesize, fill=white, inner sep=1.5pt},
}

\title{\textbf{Email Triage Agent}\\[0.3em]
\large Agent AI do automatycznej obsługi zgłoszeń klientów\\
z weryfikacją przez człowieka (human-in-the-loop)}
\author{Michał Pawlik}
\date{\today}

\begin{document}

\maketitle

\tableofcontents
\newpage

% =====================================================================
\section{Cel projektu}

Firmy dostają codziennie dziesiątki lub setki maili od klientów. Duża część
z nich to powtarzalne pytania (np. \emph{,,ile trwa dostawa?''}), na które
odpowiedź jest zawsze podobna. Część jest jednak wrażliwa (reklamacja, zwrot
pieniędzy, sprawa prawna) i na nią nie powinien odpowiadać automat bez
nadzoru.

\textbf{Email Triage Agent} to program, który dla każdego przychodzącego
zgłoszenia:

\begin{enumerate}
  \item \textbf{klasyfikuje} je, czyli określa jego kategorię i to, jak
        pewna jest ta klasyfikacja,
  \item \textbf{wyszukuje} w bazie wiedzy firmy fakty potrzebne do odpowiedzi
        (technika RAG),
  \item \textbf{pisze szkic odpowiedzi} na podstawie znalezionych faktów,
  \item \textbf{decyduje}, czy odpowiedź może zostać wysłana automatycznie,
        czy musi ją najpierw zatwierdzić człowiek.
\end{enumerate}

Spam i phishing są rozpoznawane na samym początku i odrzucane bez
generowania jakiejkolwiek odpowiedzi.

Najważniejsza idea projektu: \textbf{automat obsługuje tylko to, co jest
bezpieczne i jednoznaczne, a wszystko inne przekazuje człowiekowi}.

% =====================================================================
\section{Wykorzystane technologie}

\begin{center}
\renewcommand{\arraystretch}{1.25}
\begin{tabularx}{\textwidth}{@{}l X@{}}
\toprule
\textbf{Technologia} & \textbf{Do czego służy w projekcie} \\
\midrule
Python 3 & Język, w którym napisany jest cały program. \\
LangGraph & Biblioteka do budowy agentów jako \emph{grafu kroków}. Definiuje
  kolejność kroków, warunki rozgałęzień oraz pauzę na decyzję człowieka. \\
LangChain + Claude (Anthropic) & Wywołania modeli językowych. Claude Haiku
  klasyfikuje zgłoszenia, Claude Sonnet pisze odpowiedzi. \\
ChromaDB & Wektorowa baza danych, w której przechowywana jest baza wiedzy
  do wyszukiwania semantycznego. \\
sentence-transformers & Lokalny model zamieniający tekst na wektor liczb
  (embedding); obsługuje język polski. \\
SQLite & Trwały zapis stanu każdego zgłoszenia (tzw. checkpointer). \\
Streamlit & Prosty interfejs webowy dla człowieka zatwierdzającego odpowiedzi. \\
Gmail API & Opcjonalne źródło prawdziwych maili (tylko odczyt). \\
pytest & Testy automatyczne. \\
\bottomrule
\end{tabularx}
\end{center}

% =====================================================================
\section{Architektura: agent jako graf}

\subsection{Czym jest graf w LangGraph}

Program nie jest jednym dużym skryptem, tylko \textbf{grafem}, w którym:

\begin{itemize}
  \item \textbf{węzeł} (ang. \emph{node}) to jeden krok przetwarzania,
        zaimplementowany jako zwykła funkcja Pythona,
  \item \textbf{krawędź} (ang. \emph{edge}) mówi, który krok wykonuje się
        po którym,
  \item \textbf{krawędź warunkowa} wybiera następny krok na podstawie
        dotychczasowych wyników (np. ,,jeśli spam, to idź do \texttt{skip}'').
\end{itemize}

Wszystkie węzły pracują na jednym wspólnym obiekcie, czyli \textbf{stanie}
zgłoszenia (rozdział~\ref{sec:stan}). Każdy węzeł dostaje stan, dopisuje do
niego swoje wyniki i oddaje go dalej.

\subsection{Schemat grafu}

Rysunek~\ref{fig:graf} pokazuje pełny przepływ zgłoszenia. Węzły na
pomarańczowo wywołują model językowy (kosztują), węzły na niebiesko to zwykły
kod Pythona.

\begin{figure}[ht]
\centering
\begin{tikzpicture}[node distance=0.75cm]
  \node[start]                               (in)   {Nowe zgłoszenie};
  \node[llm, below=of in]                    (cls)  {\textbf{classify}\\ \footnotesize klasyfikacja};
  \node[wezel, below=of cls]                 (ret)  {\textbf{retrieve}\\ \footnotesize wyszukanie w bazie wiedzy};
  \node[llm, below=of ret]                   (gen)  {\textbf{generate}\\ \footnotesize szkic odpowiedzi};
  \node[wezel, below=of gen]                 (rou)  {\textbf{route}\\ \footnotesize decyzja};
  \node[wezel, below left=1.0cm and 0.6cm of rou]  (snd)  {\textbf{send}\\ \footnotesize status: sent};
  \node[wezel, below right=1.0cm and 0.6cm of rou] (hr)   {\textbf{human\_review}\\ \footnotesize pauza};
  \node[wezel, left=3.0cm of ret]            (skp)  {\textbf{skip}\\ \footnotesize status: ignored};
  \node[koniec, below=3.2cm of rou]          (end)  {KONIEC};
  \node[zew, right=1.0cm of hr]              (ui)   {Człowiek\\w Streamlit};

  \draw[strz] (in) -- (cls);
  \draw[strz] (cls) -- node[etykieta, right] {pozostałe} (ret);
  \draw[strz] (cls) -| node[etykieta, pos=0.25, above] {spam} (skp);
  \draw[strz] (ret) -- (gen);
  \draw[strz] (gen) -- (rou);
  \draw[strz] (rou) -- node[etykieta, left, pos=0.4] {auto} (snd);
  \draw[strz] (rou) -- node[etykieta, right, pos=0.4] {human\_review} (hr);
  \draw[strz] (snd) |- (end);
  \draw[strz] (hr)  |- node[etykieta, pos=0.25, right] {sent / rejected} (end);
  \draw[strz] (skp) |- (end);
  \draw[strz, dashed] ([yshift=2mm]hr.east) -- node[etykieta, above] {\footnotesize interrupt()} ([yshift=2mm]ui.west);
  \draw[strz, dashed] ([yshift=-2mm]ui.west) -- node[etykieta, below] {\footnotesize resume} ([yshift=-2mm]hr.east);
\end{tikzpicture}
\caption{Graf agenta. Pomarańczowe węzły wywołują model Claude.}
\label{fig:graf}
\end{figure}

Graf budowany jest w pliku \plik{src/graph.py}. Fragment definicji:

\begin{lstlisting}
graph.set_entry_point("classify")
graph.add_conditional_edges("classify", classify_condition,
                            {"spam": "skip", "process": "retrieve"})
graph.add_edge("retrieve", "generate")
graph.add_edge("generate", "route")
graph.add_conditional_edges("route", route_condition,
                            {"auto": "send", "human_review": "human_review"})
\end{lstlisting}

% =====================================================================
\section{Stan zgłoszenia}
\label{sec:stan}

Stan to słownik Pythona (\texttt{TypedDict}) zdefiniowany w
\plik{src/state.py}. Na początku zawiera tylko dane maila, a z każdym
krokiem dochodzą nowe pola:

\begin{center}
\renewcommand{\arraystretch}{1.2}
\begin{tabularx}{\textwidth}{@{}>{\raggedright\arraybackslash}p{4.2cm} >{\raggedright\arraybackslash}p{2.6cm} X@{}}
\toprule
\textbf{Pole} & \textbf{Ustawia} & \textbf{Znaczenie} \\
\midrule
\texttt{ticket\_id}, \texttt{sender}, \texttt{subject}, \texttt{body}
  & wejście & identyfikator, nadawca, temat i treść maila \\
\texttt{category} & classify & kategoria zgłoszenia \\
\texttt{confidence} & classify & pewność klasyfikacji, liczba od 0 do 1 \\
\texttt{retrieved\_context} & retrieve & fragmenty bazy wiedzy pasujące do maila \\
\texttt{draft\_reply} & generate & szkic odpowiedzi \\
\texttt{route\_decision} & route & \texttt{auto} albo \texttt{human\_review} \\
\texttt{status} & send / skip / human\_review & \texttt{sent}, \texttt{rejected}
  albo \texttt{ignored} \\
\texttt{audit\_log} & każdy węzeł & lista wpisów opisujących decyzje agenta \\
\bottomrule
\end{tabularx}
\end{center}

\textbf{Log audytowy} (\texttt{audit\_log}) jest ważny z punktu widzenia
bezpieczeństwa: dla każdego zgłoszenia da się potem odtworzyć, \emph{dlaczego}
agent podjął taką, a nie inną decyzję.

% =====================================================================
\section{Kroki agenta po kolei}

Każdy krok to osobny plik w katalogu \plik{src/nodes/}. Każdy węzeł ma tę samą
budowę: przyjmuje stan, dopisuje swoje pola i wpis do logu, zwraca nowy stan.

\subsection{Krok 1: classify (klasyfikacja)}

\textbf{Plik:} \plik{src/nodes/classify.py}\quad
\textbf{Model:} Claude Haiku (tani i szybki), \texttt{temperature=0}

Model dostaje temat i treść maila oraz polecenie, by odpowiedział
\emph{wyłącznie} JSON-em z kategorią i pewnością, np.:

\begin{lstlisting}
{"category": "reklamacja", "confidence": 0.93}
\end{lstlisting}

Dozwolone kategorie (zdefiniowane w \plik{src/config.py}):
\texttt{pytanie\_produktowe}, \texttt{reklamacja}, \texttt{zwrot\_pieniedzy},
\texttt{sprawa\_prawna}, \texttt{rezygnacja\_z\_uslugi}, \texttt{spam},
\texttt{inne}.

\textbf{Dlaczego Haiku i temperatura 0?} Klasyfikacja to proste zadanie
z zamkniętą listą odpowiedzi, więc nie potrzeba najdroższego modelu.
Temperatura 0 sprawia, że model odpowiada możliwie powtarzalnie.

\textbf{Zabezpieczenia przed błędną odpowiedzią modelu.} Model językowy nie
zawsze trzyma się instrukcji, więc odpowiedź jest sprawdzana, zanim zostanie
użyta:

\begin{itemize}
  \item z odpowiedzi wycinany jest fragment od pierwszego \texttt{\{} do
        ostatniego \texttt{\}}, więc tekst dookoła JSON-a (np. znaczniki
        Markdown) nie przeszkadza,
  \item kategoria jest normalizowana: \texttt{Reklamacja} staje się
        \texttt{reklamacja}, a \texttt{Zwrot pieniedzy} staje się
        \texttt{zwrot\_pieniedzy},
  \item pewność, która nie jest liczbą, zamieniana jest na 0, a wartość
        spoza zakresu 0--1 jest do niego przycinana,
  \item jeśli odpowiedzi \textbf{nie da się odczytać} albo kategoria jest
        \textbf{spoza listy}, agent nie przerywa pracy. Ustawia kategorię
        \texttt{inne} i pewność 0, więc zgłoszenie na pewno trafi do
        człowieka. W logu zapisywana jest przyczyna.
\end{itemize}

\textbf{Odcięcie spamu.} Zaraz po klasyfikacji funkcja
\texttt{classify\_condition} sprawdza, czy kategoria to \texttt{spam}. Jeśli
tak, zgłoszenie od razu idzie do węzła \texttt{skip}, z pominięciem reszty
grafu. Powody są dwa: odpowiedź na phishing potwierdziłaby oszustowi, że
skrzynka jest aktywna, a generowanie odpowiedzi na spam to niepotrzebny koszt.

\subsection{Krok 2: retrieve (wyszukiwanie w bazie wiedzy, RAG)}

\textbf{Pliki:} \plik{src/nodes/retrieve.py},
\plik{src/tools/knowledge_base.py}\quad
\textbf{Dane:} \plik{data/knowledge_base.json} (30 krótkich faktów o firmie)

\textbf{RAG} (\emph{Retrieval-Augmented Generation}) oznacza, że zanim model
napisze odpowiedź, program wyszukuje dla niego właściwe informacje. Dzięki
temu model opiera się na faktach z firmy, a nie na swojej ogólnej wiedzy.

Jak działa wyszukiwanie:

\begin{enumerate}
  \item Przy starcie programu każdy wpis bazy wiedzy zamieniany jest na
        \textbf{embedding}, czyli wektor liczb opisujący \emph{znaczenie}
        tekstu. Wektory zapisywane są w bazie Chroma (\plik{data/chroma/}).
  \item Temat i treść maila też zamieniane są na wektor.
  \item Chroma zwraca 3 wpisy, których wektory są najbliżej wektora maila,
        czyli wpisy o najbardziej zbliżonym znaczeniu.
\end{enumerate}

Przykład: mail \emph{,,Kiedy dotrze moje zamówienie?''} znajdzie wpis
\emph{,,Czas dostawy standardowej wynosi 3--5 dni roboczych\ldots''}, choć
nie mają prawie żadnych wspólnych słów. Tym różni się wyszukiwanie
semantyczne od zwykłego szukania słów kluczowych.

Użyty model embeddingów to \texttt{paraphrase-multilingual-MiniLM-L12-v2}.
Działa lokalnie, bez klucza API, i obsługuje język polski. Domyślny model
Chroma jest anglojęzyczny i w testach źle dopasowywał polskie pytania.

\subsection{Krok 3: generate (pisanie odpowiedzi)}

\textbf{Plik:} \plik{src/nodes/generate.py}\quad
\textbf{Model:} Claude Sonnet (lepsza jakość języka)

Model dostaje kategorię, treść maila i fragmenty z bazy wiedzy. Prompt
zawiera dwie kluczowe zasady ograniczające tzw. \textbf{halucynacje}
(zmyślanie faktów przez model):

\begin{itemize}
  \item używaj \emph{wyłącznie} faktów z podanego kontekstu i nie dodawaj
        warunków ani terminów, których w nim nie ma,
  \item jeśli kontekst nie odpowiada na pytanie klienta, napisz to wprost,
        zamiast zgadywać.
\end{itemize}

Szczegół techniczny: model Sonnet może najpierw ,,myśleć'' (tzw. adaptive
thinking) i zwrócić odpowiedź jako listę bloków (blok myślenia i blok
tekstu). Funkcja \texttt{\_extract\_text()} wybiera z niej tylko właściwy
tekst odpowiedzi.

\subsection{Krok 4: route (decyzja: automat czy człowiek)}

\textbf{Plik:} \plik{src/nodes/route.py}

To najważniejszy krok z punktu widzenia bezpieczeństwa. Nie korzysta z modelu
językowego, tylko z prostych, przewidywalnych reguł:

\begin{lstlisting}
is_unknown        = category not in CATEGORIES
is_sensitive      = category in SENSITIVE_CATEGORIES
is_low_confidence = confidence < CONFIDENCE_THRESHOLD   # 0.85

decision = "human_review" if (is_unknown or is_sensitive
                              or is_low_confidence) else "auto"
\end{lstlisting}

Kategorie wrażliwe, które \textbf{zawsze} trafiają do człowieka, niezależnie
od pewności: \texttt{reklamacja}, \texttt{zwrot\_pieniedzy},
\texttt{sprawa\_prawna}, \texttt{rezygnacja\_z\_uslugi}.

\begin{center}
\renewcommand{\arraystretch}{1.2}
\begin{tabular}{@{}l c l@{}}
\toprule
\textbf{Przykład} & \textbf{Pewność} & \textbf{Decyzja} \\
\midrule
pytanie\_produktowe & 0.95 & auto (wysłane automatycznie) \\
pytanie\_produktowe & 0.60 & human\_review (niska pewność) \\
reklamacja          & 0.99 & human\_review (kategoria wrażliwa) \\
nieznana kategoria  & --   & human\_review (nieznana kategoria) \\
\bottomrule
\end{tabular}
\end{center}

Próg pewności można zmienić bez ruszania kodu, przez zmienną
\texttt{CONFIDENCE\_THRESHOLD} w pliku \plik{.env}.

\subsection{Krok 5: zakończenie (send, skip, human\_review)}

\textbf{Plik:} \plik{src/nodes/send.py}

\begin{itemize}
  \item \textbf{send}: ustawia status \texttt{sent}. Na tym etapie projektu
        jest to symulacja, a podłączenie prawdziwej wysyłki maili to
        planowany kolejny krok.
  \item \textbf{skip}: dla spamu ustawia status \texttt{ignored}. Odpowiedź
        nie jest nawet generowana.
  \item \textbf{human\_review}: zatrzymuje agenta i czeka na człowieka.
        Opisane dokładnie w następnym rozdziale.
\end{itemize}

% =====================================================================
\section{Human-in-the-loop: jak agent czeka na człowieka}

\subsection{Pauza i wznowienie}

Węzeł \texttt{human\_review} wywołuje funkcję \texttt{interrupt()}
z LangGraph. Działa ona jak ,,pauza'': wykonanie grafu zatrzymuje się w tym
miejscu, a dane do przeglądu (nadawca, temat, kategoria, szkic odpowiedzi)
zostają przekazane na zewnątrz.

\begin{lstlisting}
decision = interrupt({
    "ticket_id": state["ticket_id"],
    "category": state.get("category"),
    "draft_reply": state.get("draft_reply"),
    ...
})
# tutaj program wraca dopiero po decyzji człowieka
\end{lstlisting}

Człowiek w interfejsie Streamlit czyta szkic, może go poprawić, i klika
,,Zaakceptuj'' albo ,,Odrzuć''. Interfejs wznawia wtedy graf specjalnym
poleceniem:

\begin{lstlisting}
agent.invoke(Command(resume={"approved": True,
                             "edited_reply": poprawiony_tekst}), config)
\end{lstlisting}

Funkcja \texttt{interrupt()} ,,zwraca'' wtedy tę decyzję, a węzeł ustawia
status \texttt{sent} (zaakceptowano) albo \texttt{rejected} (odrzucono).

\subsection{Checkpointer: pamięć między programami}

Żeby pauza mogła trwać dowolnie długo, a wznowienie nastąpić \emph{w innym
programie}, stan każdego zgłoszenia jest zapisywany na dysku w bazie SQLite
(\plik{data/checkpoints.sqlite}). Odpowiada za to tzw. \textbf{checkpointer}
(\plik{src/checkpointer.py}).

Każde zgłoszenie ma własny wątek (\texttt{thread\_id}) równy jego
identyfikatorowi (\texttt{ticket\_id}). Dzięki temu można np.:

\begin{enumerate}
  \item uruchomić \plik{app.py}, który przetworzy maile i zostawi reklamację
        ,,w pauzie'',
  \item zamknąć program,
  \item później uruchomić \plik{streamlit\_app.py}, który odczyta z bazy to
        samo zgłoszenie i pozwoli je dokończyć.
\end{enumerate}

\begin{figure}[ht]
\centering
\begin{tikzpicture}[node distance=1.2cm]
  \node[wezel, minimum width=3.2cm] (app) {\plik{app.py}\\ \footnotesize lub \plik{gmail\_app.py}};
  \node[koniec, cylinder, shape border rotate=90, aspect=0.25,
        minimum width=3.2cm, minimum height=1.4cm, right=1.6cm of app] (db)
        {\plik{checkpoints.sqlite}};
  \node[wezel, minimum width=3.2cm, right=1.6cm of db] (st) {\plik{streamlit\_app.py}};
  \node[wezel, minimum width=3.2cm, below=1.0cm of db] (met) {\plik{metrics.py}};

  \draw[strz] (app) -- node[etykieta, above] {zapis} (db);
  \draw[strz, {Stealth[length=2.2mm]}-{Stealth[length=2.2mm]}] (db) -- node[etykieta, above] {odczyt / wznowienie} (st);
  \draw[strz] (db) -- node[etykieta, right] {odczyt} (met);
\end{tikzpicture}
\caption{Programy współdzielą stan zgłoszeń przez jedną bazę SQLite.}
\end{figure}

% =====================================================================
\section{Programy do uruchomienia}

\subsection{Źródła zgłoszeń}

Oba programy wejściowe korzystają z tej samej pętli (\plik{src/cli.py}),
która przepuszcza każde zgłoszenie przez graf i wypisuje wynik w terminalu.

\begin{itemize}
  \item \plik{app.py}: wersja demonstracyjna. Czyta 3 przykładowe
        zgłoszenia z \plik{data/sample_tickets.json}: pytanie o dostawę,
        reklamację uszkodzonej przesyłki i pytanie o reset hasła.
  \item \plik{gmail\_app.py}: czyta prawdziwe, nieprzeczytane maile ze
        skrzynki Gmail (\plik{src/tools/gmail\_client.py}). Logowanie odbywa
        się przez OAuth z uprawnieniem \textbf{tylko do odczytu}
        (\texttt{gmail.readonly}), więc agent fizycznie nie może niczego
        wysłać, usunąć ani oznaczyć jako przeczytane.
\end{itemize}

\subsection{Interfejs przeglądu (Streamlit)}

\plik{streamlit\_app.py} to prosta strona w przeglądarce. Pokazuje listę
zgłoszeń, przy czym czekające na decyzję oznaczone są jako ,,[czeka]'' i
wyświetlane na górze. Dla wybranego zgłoszenia widać treść maila, kategorię,
pewność, edytowalny szkic odpowiedzi, przyciski decyzji oraz log audytowy.

\subsection{Metryki}

\plik{metrics.py} odczytuje z bazy SQLite wszystkie dotąd przetworzone
zgłoszenia i podaje statystyki: jaki procent obsłużono w pełni automatycznie,
ile trafiło do człowieka, ile odrzucono jako spam, rozkład kategorii oraz
średnią pewność klasyfikacji.

\subsection{Polecenia}

\begin{lstlisting}[language=bash]
pip install -r requirements.txt     # instalacja bibliotek
# w pliku .env ustawić ANTHROPIC_API_KEY=...

python app.py                       # agent na przykładowych zgłoszeniach
streamlit run streamlit_app.py      # panel akceptacji w przeglądarce
python metrics.py                   # statystyki
pytest                              # testy
\end{lstlisting}

% =====================================================================
\section{Przykładowy przebieg}

Poniżej log audytowy dwóch przykładowych zgłoszeń (wartości pewności są
przykładowe, w praktyce zależą od odpowiedzi modelu).

\textbf{T-1001, ,,Ile trwa dostawa?''}, czyli obsługa w pełni automatyczna:

\begin{lstlisting}[language={}]
[classify] kategoria=pytanie_produktowe pewność=0.95
[retrieve] znaleziono 3 fragmentów kontekstu
[generate] wygenerowano szkic odpowiedzi
[route] decyzja=auto (wysoka pewność, kategoria bezpieczna)
[send] wysłano odpowiedź do klient1@example.com
\end{lstlisting}

\textbf{T-1002, ,,Uszkodzona przesyłka''}, czyli przegląd przez człowieka:

\begin{lstlisting}[language={}]
[classify] kategoria=reklamacja pewność=0.95
[retrieve] znaleziono 3 fragmentów kontekstu
[generate] wygenerowano szkic odpowiedzi
[route] decyzja=human_review (kategoria wrażliwa)
      ... pauza, decyzja człowieka w Streamlit ...
[human_review] zaakceptowano przez człowieka - wysłano do klient2@example.com
\end{lstlisting}

% =====================================================================
\section{Konfiguracja}

Ustawienia wczytywane są z pliku \plik{.env} (\plik{src/config.py}), więc
zachowanie agenta można zmieniać bez modyfikowania kodu.

\begin{center}
\renewcommand{\arraystretch}{1.2}
\begin{tabularx}{\textwidth}{@{}l l X@{}}
\toprule
\textbf{Zmienna} & \textbf{Domyślnie} & \textbf{Znaczenie} \\
\midrule
\texttt{ANTHROPIC\_API\_KEY} & -- & klucz API do modeli Claude (wymagany) \\
\texttt{CLASSIFY\_MODEL\_NAME} & \texttt{claude-haiku-4-5} & model do klasyfikacji \\
\texttt{GENERATE\_MODEL\_NAME} & \texttt{claude-sonnet-5} & model do pisania odpowiedzi \\
\texttt{CONFIDENCE\_THRESHOLD} & \texttt{0.85} & próg pewności dla automatu \\
\texttt{GMAIL\_QUERY} & \texttt{is:unread in:inbox} & które maile pobierać z Gmaila \\
\texttt{GMAIL\_MAX\_RESULTS} & \texttt{10} & ile maili pobrać naraz \\
\texttt{LANGSMITH\_TRACING} & wyłączone & opcjonalny podgląd każdego kroku w LangSmith \\
\bottomrule
\end{tabularx}
\end{center}

Pliki z kluczami (\plik{.env}, token Gmaila, \plik{client\_secret*.json})
są wykluczone z repozytorium Git (\plik{.gitignore}).

% =====================================================================
\section{Testy}

\begin{itemize}
  \item \plik{tests/test\_classify.py}: testy jednostkowe bez wywołań
        modelu, szybkie i powtarzalne. Sprawdzają odczytywanie JSON-a
        (poprawnego, otoczonego tekstem i całkowicie błędnego), normalizację
        kategorii i pewności oraz to, że nieznana kategoria lub błędna
        odpowiedź modelu zawsze kończy się przeglądem przez człowieka.
  \item \plik{tests/test\_graph.py}: testy całego grafu na prawdziwym
        modelu Claude. Bezpieczne pytanie przechodzi do końca; reklamacja
        zatrzymuje się na \texttt{human\_review} i może zostać odrzucona;
        phishing kończy się statusem \texttt{ignored} bez żadnej odpowiedzi.
\end{itemize}

% =====================================================================
\section{Mechanizmy bezpieczeństwa: podsumowanie}

\begin{enumerate}
  \item \textbf{Kategorie wrażliwe zawsze trafiają do człowieka}, nawet przy
        100\% pewności modelu.
  \item \textbf{Niska pewność trafia do człowieka.} Automat działa tylko
        w jednoznacznych przypadkach.
  \item \textbf{Nieznana kategoria lub nieczytelna odpowiedź modelu trafia
        do człowieka.} W razie wątpliwości decyduje człowiek, a program się
        nie zawiesza.
  \item \textbf{Spam i phishing nie dostają odpowiedzi.}
  \item \textbf{Odpowiedzi opierają się na bazie wiedzy (RAG)}, a prompt
        zabrania dopisywania faktów spoza niej.
  \item \textbf{Decyzja o wysyłce jest podejmowana regułami w kodzie}, a nie
        przez model, więc jest przewidywalna i łatwa do sprawdzenia.
  \item \textbf{Log audytowy} dla każdego zgłoszenia.
  \item \textbf{Gmail tylko do odczytu}, więc agent nie może niczego zmienić
        na koncie.
\end{enumerate}

% =====================================================================
\section{Ograniczenia i możliwy dalszy rozwój}

\begin{itemize}
  \item \textbf{Wysyłka jest symulowana.} Węzeł \texttt{send} tylko ustawia
        status. Kolejnym krokiem jest podłączenie wysyłki przez Gmail API.
  \item \textbf{Pewność to samoocena modelu.} Liczbę \texttt{confidence}
        model podaje sam, więc nie jest to prawdopodobieństwo w sensie
        statystycznym. Próg warto dobrać na podstawie większej liczby
        rzeczywistych maili.
  \item \textbf{Mała baza wiedzy.} 30 wpisów wystarcza do demonstracji;
        w realnym zastosowaniu bazę można zasilić np. dokumentami firmy.
  \item \textbf{Testy end-to-end korzystają z prawdziwego modelu}, więc
        kosztują i mogą dawać nieco różne wyniki między uruchomieniami.
\end{itemize}

% =====================================================================
\section{Struktura plików}

\begin{lstlisting}[language={}]
email-triage-agent/
  app.py                 # uruchomienie na przykładowych zgłoszeniach
  gmail_app.py           # uruchomienie na mailach z Gmaila
  streamlit_app.py       # panel akceptacji dla człowieka
  metrics.py             # statystyki
  requirements.txt       # lista bibliotek
  .env.example           # wzór pliku konfiguracyjnego
  data/
    sample_tickets.json  # 3 przykładowe zgłoszenia
    knowledge_base.json  # baza wiedzy (30 faktów)
  src/
    state.py             # definicja stanu zgłoszenia
    graph.py             # budowa grafu agenta
    config.py            # konfiguracja, progi, kategorie
    checkpointer.py      # zapis stanu w SQLite
    cli.py               # wspólna pętla przetwarzania
    nodes/
      classify.py        # krok 1: klasyfikacja
      retrieve.py        # krok 2: wyszukiwanie (RAG)
      generate.py        # krok 3: szkic odpowiedzi
      route.py           # krok 4: decyzja auto / człowiek
      send.py            # krok 5: send, skip, human_review
    tools/
      knowledge_base.py  # baza wektorowa Chroma
      gmail_client.py    # odczyt maili z Gmaila
  tests/
    test_classify.py     # testy jednostkowe
    test_graph.py        # testy całego agenta
\end{lstlisting}

\end{document}
